\documentclass{article}

% if you need to pass options to natbib, use, e.g.:
%     \PassOptionsToPackage{numbers, compress}{natbib}
% before loading neurips_2025

% The authors should use one of these tracks.
% Before accepting by the NeurIPS conference, select one of the options below.
% 0. "default" for submission
 \usepackage{neurips_2025}
% the "default" option is equal to the "main" option, which is used for the Main Track with double-blind reviewing.
% 1. "main" option is used for the Main Track
%  \usepackage[main]{neurips_2025}
% 2. "position" option is used for the Position Paper Track
%  \usepackage[position]{neurips_2025}
% 3. "dandb" option is used for the Datasets & Benchmarks Track
 % \usepackage[dandb]{neurips_2025}
% 4. "creativeai" option is used for the Creative AI Track
%  \usepackage[creativeai]{neurips_2025}
% 5. "sglblindworkshop" option is used for the Workshop with single-blind reviewing
 % \usepackage[sglblindworkshop]{neurips_2025}
% 6. "dblblindworkshop" option is used for the Workshop with double-blind reviewing
%  \usepackage[dblblindworkshop]{neurips_2025}

% After being accepted, the authors should add "final" behind the track to compile a camera-ready version.
% 1. Main Track
 % \usepackage[main, final]{neurips_2025}
% 2. Position Paper Track
%  \usepackage[position, final]{neurips_2025}
% 3. Datasets & Benchmarks Track
 % \usepackage[dandb, final]{neurips_2025}
% 4. Creative AI Track
%  \usepackage[creativeai, final]{neurips_2025}
% 5. Workshop with single-blind reviewing
%  \usepackage[sglblindworkshop, final]{neurips_2025}
% 6. Workshop with double-blind reviewing
%  \usepackage[dblblindworkshop, final]{neurips_2025}
% Note. For the workshop paper template, both \title{} and \workshoptitle{} are required, with the former indicating the paper title shown in the title and the latter indicating the workshop title displayed in the footnote.
% For workshops (5., 6.), the authors should add the name of the workshop, "\workshoptitle" command is used to set the workshop title.
% \workshoptitle{WORKSHOP TITLE}

% "preprint" option is used for arXiv or other preprint submissions
 % \usepackage[preprint]{neurips_2025}

% to avoid loading the natbib package, add option nonatbib:
%    \usepackage[nonatbib]{neurips_2025}

\usepackage[utf8]{inputenc} % allow utf-8 input
\usepackage[T1]{fontenc}    % use 8-bit T1 fonts
\usepackage{hyperref}       % hyperlinks
\usepackage{url}            % simple URL typesetting
\usepackage{booktabs}       % professional-quality tables
\usepackage{amsfonts}       % blackboard math symbols
\usepackage{nicefrac}       % compact symbols for 1/2, etc.
\usepackage{microtype}      % microtypography
\usepackage{xcolor}         % colors

% Note. For the workshop paper template, both \title{} and \workshoptitle{} are required, with the former indicating the paper title shown in the title and the latter indicating the workshop title displayed in the footnote. 
\title{Parameter-Agnostic LMO-Based Ringmaster ASGD}


% The \author macro works with any number of authors. There are two commands
% used to separate the names and addresses of multiple authors: \And and \AND.
%
% Using \And between authors leaves it to LaTeX to determine where to break the
% lines. Using \AND forces a line break at that point. So, if LaTeX puts 3 of 4
% authors names on the first line, and the last on the second line, try using
% \AND instead of \And before the third author name.


\author{%
  David S.~Hippocampus\thanks{Use footnote for providing further information
    about author (webpage, alternative address)---\emph{not} for acknowledging
    funding agencies.} \\
  Department of Computer Science\\
  Cranberry-Lemon University\\
  Pittsburgh, PA 15213 \\
  \texttt{hippo@cs.cranberry-lemon.edu} \\
  % examples of more authors
  % \And
  % Coauthor \\
  % Affiliation \\
  % Address \\
  % \texttt{email} \\
  % \AND
  % Coauthor \\
  % Affiliation \\
  % Address \\
  % \texttt{email} \\
  % \And
  % Coauthor \\
  % Affiliation \\
  % Address \\
  % \texttt{email} \\
  % \And
  % Coauthor \\
  % Affiliation \\
  % Address \\
  % \texttt{email} \\
}

\usepackage{graphicx}
\usepackage{apptools}
\usepackage[flushleft]{threeparttable}
\usepackage{array,booktabs,makecell}
\usepackage{multirow}
\usepackage{nicefrac}
\usepackage{amsthm}
\usepackage{mathtools}
\usepackage{amsmath,amsfonts,bm,amssymb}
\usepackage{dsfont}
\usepackage{wrapfig}
\usepackage{caption}
% \usepackage{subcaption}
\usepackage{siunitx}
\usepackage{thm-restate}
\usepackage{nccmath}
\usepackage{empheq}
\usepackage{bbm}
\usepackage{tabularx}

% algorithmic and algorithm already loaded by icml2026.sty
\usepackage{algorithmic}
\usepackage{algorithm}
\usepackage{filecontents}

\usepackage[capitalize,noabbrev]{cleveref}

% COLORS
\usepackage{color}
\usepackage{colortbl}
\definecolor{bgcolor}{rgb}{0.76,0.88,0.50}
\definecolor{bgcolor0}{rgb}{0.93,0.99,1}
\definecolor{bgcolor1}{rgb}{0.8,1,1}
\definecolor{bgcolor2}{rgb}{0.8,1,0.8}
\definecolor{bgcolor3}{rgb}{0.50,0.90,0.50}
\usepackage{tcolorbox}
\usepackage{pifont}

\definecolor{mydarkblue}{rgb}{0,0.16,0.54}
\definecolor{mydarkgreen}{RGB}{39,130,67}
\definecolor{mydarkorange}{RGB}{236,147,14}
% \definecolor{myorange}{RGB}{230,120,20}
% \definecolor{mydarkred}{RGB}{192,47,25}
\definecolor{mydarkgreen}{RGB}{0,100,0}
\definecolor{ruby}{RGB}{155,17,30}
\definecolor{chili}{RGB}{191,0,0}
\definecolor{sangria}{RGB}{146,0,10}
\definecolor{burgundy}{RGB}{128,0,32}
\definecolor{darkred}{RGB}{132,0,0} 
\definecolor{cherry}{RGB}{192,0,0} 
\definecolor{myaccent}{rgb}{0,0.45,0.45}
\definecolor{anotherdarkred}{rgb}{0.45,0,0.05}
\newcommand{\green}{\color{mydarkgreen}}
\newcommand{\orange}{\color{mydarkorange}}
\newcommand{\blue}{\color{mydarkblue}}
\newcommand{\red}{\color{cherry}}
% \newcommand{\red}{\color{anotherdarkred}}
\newcommand{\highlightcolor}{\color{myaccent}}

\newcommand{\markchanges}{}

% todonotes already loaded in paper.tex
\usepackage[textsize=tiny]{todonotes}
\newcommand{\arto}[1]{\todo[inline]{\textbf{Arto: }#1}}
\newcommand{\zhiro}[1]{\todo[inline]{\textbf{Zhirayr: }#1}}
\newcommand{\peter}[1]{\todo[inline]{\textbf{Peter: }#1}}

\usepackage{xspace}
\usepackage{relsize}
\newcommand{\dataname}[1]{{\tt\footnotesize\color{blue}#1}\xspace}


\newcommand{\ind}{\perp\!\!\!\!\perp}
\newcommand{\norm}[1]{\left\| #1 \right\|}
\newcommand{\sqnorm}[1]{\left\| #1 \right\|^2}
\newcommand{\infnorm}[1]{\left\| #1 \right\|_{\infty}}
\newcommand{\flr}[1]{\left\lfloor #1\right\rfloor} % floor
\newcommand{\ceil}[1]{\left\lceil #1\right\rceil} % ceiling
\newcommand{\lin}[1]{\left\langle #1\right\rangle} % inner product
\newcommand{\inp}[2]{\left\langle#1,#2\right\rangle} % inner product
\newcommand{\abs}[1]{\left| #1 \right|}

\newcommand{\R}{\mathbb{R}} % reals
\newcommand{\N}{\mathbb{N}} % natural numbers
\newcommand{\Z}{\mathbb{Z}} % integers
\newcommand{\E}[1]{\mathbb{E}\left[#1\right]}
\newcommand{\Med}[1]{\mathrm{Med}\left[#1\right]}
\newcommand{\Var}[1]{\mathrm{Var}\left(#1\right)}

\newcommand{\supp}{\operatorname{supp}}
\newcommand{\progg}{\mathrm{prog}}

\newcommand{\Exp}[1]{{\mathbb{E}}\left[#1\right]}
\newcommand{\ExpSub}[2]{{\mathbb{E}}_{#1}\left[#2\right]}
\newcommand{\ExpCond}[2]{{\mathbb{E}}\left[\left.#1\ \right\vert\ #2\right]}

\newcommand{\Prob}[1]{\mathbb{P}\left(#1\right)} % probability
\newcommand{\ProbCond}[2]{\mathbb{P}\left(#1\middle\vert#2\right)}

% hat
\newcommand{\hmu}{\hat{\mu}}

% confidence
\newcommand{\conf}{\mathrm{conf}}

% caligraphic
\newcommand{\cA}{\mathcal{A}}
\newcommand{\cB}{\mathcal{B}}
\newcommand{\cC}{\mathcal{C}}
\newcommand{\cD}{\mathcal{D}}
\newcommand{\cE}{\mathcal{E}}
\newcommand{\cF}{\mathcal{F}}
\newcommand{\cG}{\mathcal{G}}
\newcommand{\cH}{\mathcal{H}}
\newcommand{\cI}{\mathcal{I}}
\newcommand{\cJ}{\mathcal{J}}
\newcommand{\cK}{\mathcal{K}}
\newcommand{\cL}{\mathcal{L}}
\newcommand{\cM}{\mathcal{M}}
\newcommand{\cN}{\mathcal{N}}
\newcommand{\cO}{\mathcal{O}}
\newcommand{\cP}{\mathcal{P}}
\newcommand{\cQ}{\mathcal{Q}}
\newcommand{\cR}{\mathcal{R}}
\newcommand{\cS}{\mathcal{S}}
\newcommand{\cT}{\mathcal{T}}
\newcommand{\cU}{\mathcal{U}}
\newcommand{\cV}{\mathcal{V}}
\newcommand{\cW}{\mathcal{W}}
\newcommand{\cX}{\mathcal{X}}
\newcommand{\cY}{\mathcal{Y}}
\newcommand{\cZ}{\mathcal{Z}}

% bold letters
\newcommand{\ba}{\bm{a}}
\newcommand{\bb}{\bm{b}}
\newcommand{\bc}{\bm{c}}
\newcommand{\bd}{\bm{d}}
\newcommand{\be}{\bm{e}}
\newcommand{\bg}{\bm{g}}
\newcommand{\bh}{\bm{h}}
\newcommand{\bi}{\bm{i}}
\newcommand{\bj}{\bm{j}}
\newcommand{\bk}{\bm{k}}
\newcommand{\bl}{\bm{l}}
\newcommand{\bn}{\bm{n}}
\newcommand{\bo}{\bm{o}}
\newcommand{\bp}{\bm{p}}
\newcommand{\bq}{\bm{q}}
\newcommand{\br}{\bm{r}}
\newcommand{\bs}{\bm{s}}
\newcommand{\bt}{\bm{t}}
\newcommand{\bu}{\bm{u}}
\newcommand{\bv}{\bm{v}}
\newcommand{\bw}{\bm{w}}
\newcommand{\bx}{\bm{x}}
\newcommand{\by}{\bm{y}}
\newcommand{\bz}{\bm{z}}

% bold matrices
\newcommand{\mA}{\mathbf{A}}
\newcommand{\mB}{\mathbf{B}}
\newcommand{\mC}{\mathbf{C}}
\newcommand{\mD}{\mathbf{D}}
\newcommand{\mE}{\mathbf{E}}
\newcommand{\mF}{\mathbf{F}}
\newcommand{\mG}{\mathbf{G}}
\newcommand{\mH}{\mathbf{H}}
\newcommand{\mI}{\mathbf{I}}
\newcommand{\mJ}{\mathbf{J}}
\newcommand{\mK}{\mathbf{K}}
\newcommand{\mL}{\mathbf{L}}
\newcommand{\mM}{\mathbf{M}}
\newcommand{\mN}{\mathbf{N}}
\newcommand{\mO}{\mathbf{O}}
\newcommand{\mP}{\mathbf{P}}
\newcommand{\mQ}{\mathbf{Q}}
\newcommand{\mR}{\mathbf{R}}
\newcommand{\mS}{\mathbf{S}}
\newcommand{\mT}{\mathbf{T}}
\newcommand{\mU}{\mathbf{U}}
\newcommand{\mV}{\mathbf{V}}
\newcommand{\mW}{\mathbf{W}}
\newcommand{\mX}{\mathbf{X}}
\newcommand{\mY}{\mathbf{Y}}
\newcommand{\mZ}{\mathbf{Z}}

\newcommand{\eqdef}{\coloneqq} 
\newcommand{\Mod}[1]{\ (\mathrm{mod}\ #1)}
\newcommand{\minimize}{\mathop{\mathrm{minimize}}}


\newcommand{\tauavg}{\tau_{\mathrm{avg}}}

\newcommand{\hnabla}{\widehat{\nabla}}


\makeatletter
\newcommand{\vast}{\bBigg@{4}}

\DeclareMathOperator*{\argmin}{arg\,min}
\DeclareMathOperator*{\argmax}{arg\,max}

\def\<{\left\langle}
\def\>{\right\rangle}
\def\[{\left[}
\def\]{\right]}
\def\({\left(}
\def\){\right)}

% The old colors
% \newcommand{\crossmarkred}{{\color{mydarkred}\ding{56}}}
% \newcommand{\checkmarkgreen}{\textbf{\color{mydarkgreen}\ding{52}}}

% % No colors
% \newcommand{\checkmarkgreen}{\textbf{\ding{52}}}
% \newcommand{\crossmarkred}{\ding{56}}

% Different colors
% \definecolor{myblue}{rgb}{0.25,0.25,0.3}
% \definecolor{myorange}{rgb}{0.25,0.25,0.3}

\definecolor{myblue}{rgb}{0.3,0.25,0.2}
\definecolor{myorange}{rgb}{0.3,0.25,0.2}
\newcommand{\checkmarkgreen}{\textbf{\color{myblue}\ding{52}}}
\newcommand{\crossmarkred}{{\color{myorange}\ding{56}}}

%%%%%%%%%%%%%%%%%%%%%%%%%%%%%%%%
% Algorithm names
%%%%%%%%%%%%%%%%%%%%%%%%%%%%%%%%
% \newcommand{\algname}[1]{{\sf\footnotesize\red#1}\xspace}
% \newcommand{\algname}[1]{{\sf\red\relscale{0.90}#1}\xspace}
\newcommand{\algname}[1]{\text{\sf\relscale{0.93}#1}\xspace}
\newcommand{\algnameS}[1]{\text{\sf\scriptsize\red#1}\xspace}
% \newcommand{\algn}{{\sf\red\relscale{0.95}Ringleader ASGD}\xspace}
\newcommand{\iasgd}{{\sf\relscale{0.95}IA\textsuperscript{2}SGD}\xspace}
\newcommand{\iasgdtitle}{IA\textsuperscript{2}SGD}
\newcommand{\ringleader}{\algname{Ringleader~ASGD}}
\newcommand{\ringleadertitle}{Ringleader~ASGD}
\newcommand{\malenia}{\algname{Malenia~SGD}}
\newcommand{\maleniatitle}{Malenia~SGD}
\newcommand{\herosgd}{\algname{Hero~SGD}}
\newcommand{\herosgdtitle}{Hero~SGD}
\newcommand{\minibatch}{\algname{Minibatch~SGD}}
\newcommand{\minibatchtitle}{Minibatch~SGD}
\newcommand{\naiveminibatch}{\algname{Naive~Minibatch~SGD}}
\newcommand{\naiveminibatchtitle}{Naive~Minibatch~SGD}
\newcommand{\asyncsgdtitle}{Asynchronous~SGD}
\newcommand{\asyncsgd}{\algname{\asyncsgdtitle}}
\newcommand{\naiveasyncsgdtitle}{Naive~Asynchronous~SGD}
\newcommand{\naiveasyncsgd}{\algname{\naiveasyncsgdtitle}}
\newcommand{\sag}{\algname{SAG}}
\newcommand{\sagtitle}{SAG}
\newcommand{\saga}{\algname{SAGA}}
\newcommand{\sagatitle}{SAGA}
\newcommand{\ringmaster}{\algname{Ringmaster~ASGD}}
\newcommand{\ringmastertitle}{Ringmaster~ASGD}
\newcommand{\rennalatitle}{Rennala~SGD}
\newcommand{\rennala}{\algname{Rennala~SGD}}
\newcommand{\sgd}{\algname{SGD}}
\newcommand{\ata}{\algname{ATA}}
\newcommand{\hogwild}{\algname{HOGWILD!}}
\newcommand{\mvrtitle}{MVR}
\newcommand{\mvr}{\algname{MVR}}
\newcommand{\storm}{\algname{STORM}}
\newcommand{\rennalamvrtitle}{{\red Rennala~MVR}}
\newcommand{\rennalamvr}{\algname{\red Rennala~MVR}}


%%%%%%%%%%%%%%%%%%%%%%%%%%%%%%%%
% THEOREMS
%%%%%%%%%%%%%%%%%%%%%%%%%%%%%%%%
\usepackage{thmtools}
\usepackage{thm-restate}
\usepackage{mdframed}

\theoremstyle{plain}
\newtheorem{theorem}{Theorem}[section]
\newtheorem{proposition}[theorem]{Proposition}
\newtheorem{lemma}[theorem]{Lemma}
\newtheorem{corollary}[theorem]{Corollary}
\theoremstyle{definition}
\newtheorem{definition}[theorem]{Definition}
\newtheorem{assumption}[theorem]{Assumption}
\theoremstyle{remark}
\newtheorem{remark}[theorem]{Remark}

% \theoremstyle{definition} % make theorems body non-italic
\theoremstyle{plain} % make theorems body italic

% Box style (used for boxed versions)
\newmdenv[
  font=\normalfont\bfseries,
  linecolor=black,
  linewidth=0.8pt,
  topline=false,
  bottomline=false,
  leftline=false,
  rightline=false,
  backgroundcolor=gray!13,
  skipabove=2pt,
  skipbelow=2pt,
  innertopmargin=-5pt,
  innerbottommargin=4pt,
  innerrightmargin=4pt,
  innerleftmargin=4pt,
]{myboxed}

% Boxed versions of theorem and lemma
\newmdtheoremenv[
  font=\normalfont\bfseries,
  linecolor=black,
  linewidth=0.8pt,
  topline=false,
  bottomline=false,
  leftline=false,
  rightline=false,
  backgroundcolor=gray!13,
  skipabove=2pt,
  skipbelow=2pt,
  innertopmargin=-5pt,
  innerbottommargin=4pt,
  innerrightmargin=4pt,
  innerleftmargin=4pt,
]{boxedtheorem}[theorem]{Theorem}

\newmdtheoremenv[
  font=\normalfont\bfseries,
  linecolor=black,
  linewidth=0.8pt,
  topline=false,
  bottomline=false,
  leftline=false,
  rightline=false,
  backgroundcolor=gray!13,
  skipabove=2pt,
  skipbelow=2pt,
  innertopmargin=-5pt,
  innerbottommargin=4pt,
  innerrightmargin=4pt,
  innerleftmargin=4pt,
]{boxedlemma}[theorem]{Lemma}

\newmdtheoremenv[
  font=\normalfont\bfseries,
  linecolor=black,
  linewidth=0.8pt,
  topline=false,
  bottomline=false,
  leftline=false,
  rightline=false,
  backgroundcolor=gray!13,
  skipabove=2pt,
  skipbelow=2pt,
  innertopmargin=-5pt,
  innerbottommargin=4pt,
  innerrightmargin=4pt,
  innerleftmargin=4pt,
]{boxeddefinition}[theorem]{Definition}

\newmdtheoremenv[
  font=\normalfont\bfseries,
  linecolor=black,
  linewidth=0.8pt,
  topline=false,
  bottomline=false,
  leftline=false,
  rightline=false,
  backgroundcolor=gray!13,
  skipabove=2pt,
  skipbelow=2pt,
  innertopmargin=-5pt,
  innerbottommargin=4pt,
  innerrightmargin=4pt,
  innerleftmargin=4pt,
]{boxedassumption}[theorem]{Assumption}

% Custom environments for restating (unboxed)
\newtheorem{innercustomthm}{Theorem}
\newenvironment{restate-theorem}[1]
  {\renewcommand\theinnercustomthm{#1}\innercustomthm}
  {\endinnercustomthm}

\newtheorem{innercustomlemma}{Lemma}
\newenvironment{restate-lemma}[1]
  {\renewcommand\theinnercustomlemma{#1}\innercustomlemma}
  {\endinnercustomlemma}

\newtheorem{innercustomproposition}{Proposition}
\newenvironment{restate-proposition}[1]
  {\renewcommand\theinnercustomproposition{#1}\innercustomproposition}
  {\endinnercustomproposition}

% Boxed restated theorem and lemma
\newenvironment{restate-boxedtheorem}[1]
  {\renewcommand\theinnercustomthm{#1}\begin{myboxed}\begin{innercustomthm}}
  {\end{innercustomthm}\end{myboxed}}

\newenvironment{restate-boxedlemma}[1]
  {\renewcommand\theinnercustomlemma{#1}\begin{myboxed}\begin{innercustomlemma}}
  {\end{innercustomlemma}\end{myboxed}}

% Sketch of Proof
\newcommand*{\sketchproofname}{Sketch of Proof}
\newenvironment{sketchproof}[1][\sketchproofname]
  {\begin{proof}[#1]}
  {\end{proof}}


\begin{document}


\maketitle


\begin{abstract}
    Parameter agnostic lmo-based \ringmaster.
    We might extend it to $(L_0, L_1)$--smoothness assumption case.

    Also, we can maybe do the same stuff for \ringleader.
\end{abstract}


\section{Problem Setup}
% 
% 
We consider the nonconvex optimization problem
\begin{equation}\label{eq:problem}
  \minimize_{x \in \R^d} 
  \left\{ f(x) \coloneqq \ExpSub{\xi \sim \cD}{f(x;\xi)} \right\} ,
\end{equation}
where $f(x; \xi)$ is the loss function evaluated on a data sample $\xi$ drawn from distribution~$\cD$, and the model is parameterized by $x \in \mathbb{R}^d$ with $d$ denoting the parameter dimensionality.
% The goal is to learn a global model that minimizes the expected loss.

We consider a distributed learning setting with $n$ workers, where each worker~$i$ has access to the same data distribution~$\cD$.
This setting is common in data centers with either unified memory or uniformly partitioned data.
Following the \textit{fixed computation model} \citep{mishchenko2022asynchronous}, we formalize the heterogeneous computation times as follows:
% 
\begin{equation}\label{eq:fixed_time}
  \begin{minipage}{0.9\linewidth}
    \centering
    Each worker~$i$ requires $\tau_i$~seconds to compute one stochastic~gradient~$\nabla f(x;\xi)$.

    Without loss of generality, we assume $0 < \tau_1 \le \tau_2 \le \cdots \le \tau_n$.
  \end{minipage}
\end{equation}
\arto{make this a formal assumption}
% 
We assume instantaneous communication (zero latency) between workers and the server in both directions.
This modeling simplification has been the standard assumption in prior work analyzing the time complexity of distributed optimization \citep{mishchenko2022asynchronous, koloskova2022sharper, tyurin2023optimal, maranjyan2025ringmaster}.

\arto{Rewrite above}

We introduce the following notation: linear minimization oracle
$$
    \textrm{lmo}(y) \eqdef \argmin_{x~:~\|x\| \le 1} \inp{y}{x} .
$$
% 
\begin{definition}
    \label{def:gen_smooth}
    Let $f:\R^d\rightarrow\R$ be differentiable.
    Then $f$ is symmetric $(L_0, L_1)$-smooth if there exist $L_0,L_1 \ge 0$ such that for any $x, y \in \R^d$
    \begin{equation}
        \label{eq:gen_smooth}
        \norm{\nabla f(x) - \nabla f(y)}_* \le \sup_{\theta \in [0,1]} \( L_0 + L_1 \norm{\nabla f(\theta x +(1-\theta)y)}_* \) \norm{x-y}.
    \end{equation}
\end{definition}
% 
% 
\section{Algorithm}
% 
% 
% 
\begin{algorithm}[H]
    \caption{LMO \ringmaster (without calculation stops)}
    \label{alg:ringmaster_lmo}
    \begin{algorithmic}[1]
        \STATE \textbf{Input:} points $x^0, m^0 \in \R^d$, stepsize $\eta_k > 0,$ delay threshold $R_k \in \N$, momentum parameter $0<\alpha_k \le 1$.
        \STATE Set $k = 0$
        \STATE Workers start computing stochastic gradients at $x^0$
        \WHILE{not terminated}
            \STATE Gradient $g^k = \nabla f\big(x^{k-\delta^k}; \xi^{k-\delta^k}_{i}\big)$ arrives from worker $i$
            \highlightcolor
            \IF{$\delta^k < R_k$}
                \color{black}
                \STATE Update the momentum: $m^{k+1} = \alpha_k g^k + (1-\alpha_k)m^{k} $
                \STATE Update the model:
                    $x^{k+1} = x^{k} + \eta_k\textrm{lmo}(m^{k+1})$
                \STATE Worker $i$ begins calculating 
                    $\nabla f\big(x^{k+1}; \xi^{k+1}_{i}\big)$
                \STATE Update the iteration number $k = k + 1$
            \ELSE
                \STATE Ignore the outdated gradient 
                    $\nabla f\big(x^{k-\delta^k}; \xi^{k-\delta^k}_{i}\big)$
                \STATE Worker $i$ begins calculating 
                    $\nabla f\big(x^{k}; \xi^{k}_{i}\big)$
            \ENDIF
        \ENDWHILE
    \end{algorithmic}
\end{algorithm}
% 
% 
\begin{algorithm}[H]
    \caption{\ringmaster (with calculation stops)}
    \label{alg:ringmasternewstops}
    \begin{algorithmic}[1]
        \STATE \textbf{Input:} points $x^0, m^0 \in \R^d$, stepsize $\eta_k > 0,$ delay threshold $R_k \in \N$, momentum parameter $0<\alpha_k \le 1$.
        \STATE Set $k = 0$
        \STATE Workers start computing stochastic gradients at $x^0$
        \WHILE{not terminated}
            \STATE {\highlightcolor Stop calculating stochastic gradients with delays $\geq R$, and start computing new ones at $x^k$ instead}
            \STATE Gradient $\nabla f\big(x^{k-\delta^k}; \xi^{k-\delta^k}_{i}\big)$ arrives from worker $i$
            \STATE Update the momentum: $m^{k+1} = \alpha_k g^k + (1-\alpha_k)m^{k} $
            \STATE Update the model:
                $x^{k+1} = x^{k} + \eta_k\textrm{lmo}(m^{k+1})$
            \STATE Worker $i$ begins calculating 
                $\nabla f\big(x^{k+1}; \xi^{k+1}_{i}\big)$
            \STATE Update the iteration number $k = k + 1$
        \ENDWHILE
    \end{algorithmic}
\end{algorithm}
\arto{Should I do while $k\le K$?}
\arto{Change highlight color}
\arto{Should we keep only one? Maybe the first one - without stops and just mention that we can stop similar to \citet{maranjyan2025ringmaster}.}
% 
% 
\section{Theoretical Analysis}
% 
\subsection{Iteration Complexity}
% 
\begin{theorem}[Iteration Complexity; Proof in \Cref{proof:iteration_complexity}]\label{thm:iteration_complexity}
    Take $\alpha_0 =1$ and $\alpha_k = \nicefrac{1}{\sqrt{k}}$, $R_k = \left\lfloor\nicefrac{1}{\alpha_k}\right\rfloor$ and $\eta_k = \nicefrac{\eta}{(k+1)^{\nicefrac{3}{4}}}$ for $k \ge 1$.
    The iteration complexity is
    \begin{equation*}
        K = \cO\left(\frac{\left(\nicefrac{\Delta_0 }{\eta} + {\color{red} \rho} \sigma +  L_0\eta\right)^4 }{\varepsilon^4} \( \log \frac{1}{\varepsilon} \)^4\right).
     \end{equation*}
\end{theorem}
\arto{Can we have $R_k = \lfloor \sqrt{k+1} \rfloor$? Same for $\alpha_k$?}
\subsection{Time Complexity}
% 
We now analyze the wall-clock time complexity of \Cref{alg:ringmaster_lmo}.
The following lemma gives the time required to complete $K$ iterations when using a square-root delay threshold.
% 
\begin{lemma}[Time Complexity for Square-Root Threshold; Proof in \Cref{proof:sqrt_threshold}]\label{lemma:time_sqrt}
    Consider \Cref{alg:ringmaster_lmo} with delay threshold $R_k = \lfloor \sqrt{k+1} \rfloor$ for $k \ge 0$.
    Under the fixed computation model \eqref{eq:fixed_time}, the wall-clock time needed to complete $K$ iterations satisfies
    \begin{equation*}
        T(K) = \mathcal{O}\left( \min_{m \in [n]}  \left(\frac{1}{m} \sum_{i=1}^m \frac{1}{\tau_i}\right)^{-1} \left( \sqrt{K} + \frac{K}{m} \right) \right).
    \end{equation*}
\end{lemma}
\arto{Change $R_k$?}
% 
\begin{proof}[Proof sketch]
    Define $\mathcal{K}_r \coloneqq \{k \mid \lfloor R_k \rfloor = r\}$, namely the set of iterations for which the threshold equals $r$.
    The schedule $R_k = \lfloor \sqrt{k+1} \rfloor$ partitions the $K$ iterations into these threshold-constant blocks.
    For each block $\mathcal{K}_r$, we apply \Cref{lem:time_R}, which is exactly the fixed-threshold lemma of \citet{maranjyan2025ringmaster} restated in the appendix, and then sum over all values $r \le \lfloor \sqrt{K} \rfloor$.
    Since $|\mathcal{K}_r| = 2r+1 = \mathcal{O}(r)$, each block contributes only a constant number of fixed-threshold chunks, which yields the claim.
\end{proof}
% 
Combining \Cref{thm:iteration_complexity} with \Cref{lemma:time_sqrt} yields the total wall-clock complexity for reaching an $\varepsilon$-stationary point.
% 
\begin{theorem}[Total Wall-Clock Complexity]\label{thm:time_complexity}
    The total wall-clock time for \Cref{alg:ringmaster_lmo} to reach an $\varepsilon$-stationary point is bounded by
    \begin{equation*}
        T(K) = \widetilde{\mathcal{O}}\left( \min_{m \in [n]} \left(\frac{1}{m} \sum_{i=1}^m \frac{1}{\tau_i}\right)^{-1} \left( \frac{\Phi^2}{\varepsilon^2} + \frac{\Phi^4}{m\varepsilon^4} \right) \right),
    \end{equation*}
    where $\Phi \coloneqq \nicefrac{\Delta_0}{\eta} + \rho \sigma + L_0\eta$.
\end{theorem}
\arto{Wall clock vs time. say somewhere that we use both...}
% 
\begin{proof}
    By \Cref{thm:iteration_complexity}, the required number of iterations satisfies
    $$
        K = \widetilde{\mathcal{O}}(\Phi^4 \varepsilon^{-4}).
    $$
    Substituting this bound into \Cref{lemma:time_sqrt} gives
    \begin{equation*}
        \sqrt{K} = \widetilde{\mathcal{O}}\left(\frac{\Phi^2}{\varepsilon^2}\right) \quad \text{and} \quad \frac{K}{m} = \widetilde{\mathcal{O}}\left(\frac{\Phi^4}{m\varepsilon^4}\right).
    \end{equation*}
    Plugging these estimates into the bound of \Cref{lemma:time_sqrt} yields the claim.
\end{proof}
% 
% 
\paragraph{Discussion.}
% 
\ringmaster \citep{maranjyan2025ringmaster} has the following time complexity
\begin{align*}
    \cO\left(\min\limits_{m \in [n]} \left(\frac{1}{m} \sum\limits_{i=1}^{m} \frac{1}{\tau_{i}}\right)^{-1} \left(\frac{L \Delta}{\varepsilon^2} + \frac{\sigma^2 L \Delta}{m \varepsilon^4}\right) \right) .
\end{align*}
% 
% 
% 
%%%%%%%%%%%%%%%%%%%%%%%%%%%%%%%%%%%%%%%%%%%%%%%%%%%%%%%%%%%%%%%%%%%%%%%%%%%%%%%%%%%%%%%%%%%%%%%%%%%%%%%%%%%%%%%%%%%%%%%%
%%%%%%%%%%%%%%%%%%%%%%%%%%%%%%%%%%%%%%%%%%%%%%%%%%%%%%%%%%%%%%%%%%%%%%%%%%%%%%%%%%%%%%%%%%%%%%%%%%%%%%%%%%%%%%%%%%%%%%%%
%%%%%%%%%%%%%%%%%%%%%%%%%%%%%%%%%%%%%%%%%%%%%%%%%%%%%%%%%%%%%%%%%%%%%%%%%%%%%%%%%%%%%%%%%%%%%%%%%%%%%%%%%%%%%%%%%%%%%%%%
{
\small
\bibliographystyle{apalike}
\bibliography{bib}
}
\newpage
\appendix
%%%%%%%%%%%%%%%%%%%%%%%%%%%%%%%%%%%%%%%%%%%%%%%%%%%%%%%%%%%%%%%%%%%%%%%%%%%%%%%%%%%%%%%%%%%%%%%%%%%%%%%%%%%%%%%%%%%%%%%%
%%%%%%%%%%%%%%%%%%%%%%%%%%%%%%%%%%%%%%%%%%%%%%%%%%%%%%%%%%%%%%%%%%%%%%%%%%%%%%%%%%%%%%%%%%%%%%%%%%%%%%%%%%%%%%%%%%%%%%%%
%%%%%%%%%%%%%%%%%%%%%%%%%%%%%%%%%%%%%%%%%%%%%%%%%%%%%%%%%%%%%%%%%%%%%%%%%%%%%%%%%%%%%%%%%%%%%%%%%%%%%%%%%%%%%%%%%%%%%%%%
\section{Auxiliary Lemmas}
% 
\begin{lemma} 
    \label{lem:gen_smooth}
    Let $f$ be $(L_0, L_1)$-smooth. Then for any $x,y\in \R^d$ it holds 
    \begin{gather*}
        f(x) \le f(y) + \inp{\nabla f(y)}{x-y} +\frac{1}{2}\left(L_0 + L_1\norm{\nabla f(y)}_*\right)\exp\left(L_1\norm{x-y}\right)\norm{x-y}^2, \\
        \norm{\nabla f(x) - \nabla f(y)}_* \le \left(L_0 + L_1\norm{\nabla f(y)}_*\right)\exp\left(L_1\norm{x-y}\right)\norm{x-y} .
    \end{gather*}
\end{lemma}
\arto{References?}
\begin{lemma}
    \label{lem:descent_lemma}
    Let $f$ be $(L_0, L_1)$-smooth. Then iterates of Algorithm~\ref{alg:ringmaster_lmo} satisfy 
    \begin{equation}
        \sum_{k=0}^{K-1}\eta_k\phi_k\norm{\nabla f(x^k)}_* \leq \Delta_0 + 2\sum^{K-1}_{k=0}\eta_k\norm{\hat{e}_{k+1}}_* +\frac{L_0}{2}\sum^{K-1}_{k=0}\exp\left(L_1\eta_k\right)\eta_k^2 ~,
    \end{equation}
    where we define error term as $\hat{e}_{k+1} \eqdef m^{k+1} - \nabla f(x^k)$ and constant $\phi_k  = 1-\frac{L_1\eta_k}{2}\exp(L_1\eta_k)$. 
\end{lemma}
\begin{proof}
    By Lemma~\ref{lem:gen_smooth}, we have 
    \begin{align*}
        f(x^{k+1})
        &\le f(x^k) + \inp{\nabla f(x^k)}{x^{k+1}-x^k}\\
        &\quad +\frac{1}{2}\left(L_0 + L_1\norm{\nabla f(x^k)}_*\right)\exp\left(L_1\norm{x^{k+1}-x^k}\right)\norm{x^{k+1}-x^k}^2\\
        &\le f(x^k) + \eta_k\inp{\nabla f(x^k)}{\textrm{lmo}(m^{k+1})} +\frac{1}{2}\left(L_0 + L_1\norm{\nabla f(x^k)}_*\right)\exp\left(L_1\eta_k\right)\eta_k^2\\
        &= f(x^k) + \eta_k\inp{\nabla f(x^k) - m^{k+1}}{\textrm{lmo}(m^{k+1})} + \eta_k\inp{m^{k+1}}{\textrm{lmo}(m^{k+1})} \\
        &\quad +\frac{1}{2}\left(L_0 + L_1\norm{\nabla f(x^k)}_*\right)\exp\left(L_1\eta_k\right)\eta_k^2\\
        &\le f(x^k) + \eta_k\norm{\nabla f(x^k) - m^{k+1}}_* - \eta_k\norm{m^{k+1}}_*\\
        &\quad + \frac{1}{2}\left(L_0 + L_1\norm{\nabla f(x^k)}_*\right)\exp\left(L_1\eta_k\right)\eta_k^2\\
        &\le f(x^k) + 2\eta_k\norm{\nabla f(x^k) - m^{k+1}}_* - \eta_k\norm{\nabla f(x^k)}_*\\
        &\quad + \frac{1}{2}\left(L_0 + L_1\norm{\nabla f(x^k)}_*\right)\exp\left(L_1\eta_k\right)\eta_k^2 ~.
    \end{align*}
    For the second inequality we used the fact that $\|\mathrm{lmo}(\cdot)\| \le 1$, then for the 3rd inequality we used Holder inequality and for the last step we used triangle inequality.
\end{proof}
%  
% 
We define
$$
\hat{e}_{k+1} \eqdef m^{k+1} - \nabla f(x^k) ~.
$$
\begin{lemma}
\label{lem:error_bound}
    Let $f$ be symmetric $(L_0, L_1)$--smooth.
    If stepsize $\eta_k$ is non-increasing and momentum parameter $\alpha_k$ defined as $\alpha_0 =1$ and $\alpha_k = \frac{1}{\sqrt{k}}$ for $k \ge 1$, then  the iterates of Algorithm~\ref{alg:ringmaster_lmo} satisfy 
    \begin{align*}
        \Exp{\norm{\hat{e}_{k+1}}_*}
        &\le \sum_{t=1}^k \prod_{j=t+1}^{k}(1-\alpha_j) \left(L_0 + L_1\Exp{\norm{\nabla f(x^t)}_*}\right)\exp(L_1\eta_{t-1})\eta_{t-1}\\
            &\quad + \sum_{t=1}^k \prod_{j=t+1}^{k}(1-\alpha_j) \alpha_t (L_0 + L_1 \Exp{\norm{\nabla f(x^t)}_*})\exp\left(L_1 R_t\eta_{t-\delta_t}\right)R_t\eta_{t-\delta_t}\\
                &\quad + {\color{red} \rho} \sigma\sqrt{\sum_{t=1}^k \prod_{j=t+1}^{k}(1-\alpha_j)^2 \alpha_t^2} ~.
    \end{align*}
\end{lemma}
\begin{proof}
Using momentum update, we derive
\begin{eqnarray*}
    \hat{e}_{k+1}
    &=& m^{k+1} - \nabla f(x^k) \\
    &=& (1-\alpha_k)m^k +\alpha_k g^k - \nabla f(x^k) \\
    &=& (1-\alpha_k)(m^k- \nabla f(x^{k-1})) + (1-\alpha_k)(\nabla f(x^{k-1}) - \nabla f(x^k)) +\alpha_k g^k - \alpha_k\nabla f(x^k) \\
    &=& (1-\alpha_k)(m^k- \nabla f(x^{k-1})) + (1-\alpha_k)(\nabla f(x^{k-1}) - \nabla f(x^k)) \\
    && +\alpha_k (\nabla f(x^{k-\delta_k}) - \nabla f(x^k))  +\alpha_k(g^k - \nabla f(x^{k-\delta_k}) ~.
\end{eqnarray*}
% 
Denoting $D_k \eqdef   \nabla f(x^{k-1}) - \nabla f(x^k)$, $S_k \eqdef \nabla f(x^{k-\delta_k}) - \nabla f(x^k)$, and $e_k = g^k - \nabla f(x^{k-\delta_k})$, we obtain
\begin{eqnarray*}
    \hat{e}_{k+1}
    &=& (1-\alpha_k) \hat{e}_{k} + (1-\alpha_k) D_k + \alpha_k S_k +\alpha_k e_k \\
    &=&\prod_{t = 0}^{k}(1-\alpha_t) \hat{e}_{0} + \sum_{t=0}^k \prod_{j=t+1}^{k}(1-\alpha_j)(1-\alpha_t) D_t \\
    && + \sum_{t=0}^k \prod_{j=t+1}^{k}(1-\alpha_j) \alpha_t S_t + \sum_{t=0}^k \prod_{j=t+1}^{k}(1-\alpha_j) \alpha_t e_t ~.
\end{eqnarray*}
Setting $\alpha_k$ as $\alpha_0 = \alpha_1 =1$ and $\alpha_k = \frac{1}{\sqrt{k}}$ for $k \ge 1$, we have
% 
\begin{eqnarray*}
    \hat{e}_{k+1} &=& \sum_{t=0}^k \prod_{j=t+1}^{k}(1-\alpha_j)(1-\alpha_t) D_t  + \sum_{t=0}^k \prod_{j=t+1}^{k}(1-\alpha_j) \alpha_t S_t + \sum_{t=0}^k \prod_{j=t+1}^{k}(1-\alpha_j) \alpha_t e_t\\
    &=& \underbrace{\sum_{t=1}^k \prod_{j=t+1}^{k}(1-\alpha_j)(1-\alpha_t) D_t}_{\eqdef\cT_1} + \underbrace{\sum_{t=1}^k \prod_{j=t+1}^{k}(1-\alpha_j) \alpha_t S_t}_{\eqdef \cT_2} + \underbrace{\sum_{t=1}^k \prod_{j=t+1}^{k}(1-\alpha_j) \alpha_t e_t}_{\eqdef \cT_3}.
\end{eqnarray*}
% 
To continue the proof, we need to bound the norm of terms $\cT_1$, $\cT_2$ and $\cT_3$. By triangle inequality, we get 
\begin{eqnarray*}
    \norm{\cT_1}_* &\le& \sum_{t=1}^k \prod_{j=t+1}^{k}(1-\alpha_j)(1-\alpha_t) \norm{D_t}_*\\
    &=& \sum_{t=1}^k \prod_{j=t+1}^{k}(1-\alpha_j) \norm{ \nabla f(x^{t-1}) - \nabla f(x^t)}_*\\
    &\le& \sum_{t=1}^k \prod_{j=t+1}^{k}(1-\alpha_j) \left(L_0 + L_1\norm{\nabla f(x^t)}_*\right)\exp(L_1\eta_{t-1})\eta_{t-1} ~.
\end{eqnarray*}
We used \Cref{lem:gen_smooth} in the last inequality.
% 
Bounding $\norm{\cT_2}_*$, we obtain 
\begin{eqnarray*}
    \norm{\cT_2}_* &\le& \sum_{t=1}^k \prod_{j=t+1}^{k}(1-\alpha_j) \alpha_t \norm{S_t}_*\\
    &=& \sum_{t=1}^k \prod_{j=t+1}^{k}(1-\alpha_j) \alpha_t \norm{\nabla f(x^{t-\delta_t}) - \nabla f(x^t)}_*\\
    &\le& \sum_{t=1}^k \prod_{j=t+1}^{k}(1-\alpha_j) \alpha_t \(L_0 + L_1 \norm{\nabla f(x^t)}_*\)\exp\left(L_1 \norm{x^t - x^{t-\delta_t}}\right)\norm{x^t - x^{t-\delta_t}}\\
    &\le& \sum_{t=1}^k \prod_{j=t+1}^{k}(1-\alpha_j) \alpha_t \(L_0 + L_1 \norm{\nabla f(x^t)}_*\)\\
    &&\quad \times \exp\left(L_1 \norm{\sum^{t-1}_{\tau = t-\delta_t}\eta_{\tau}\textrm{lmo}(m^{\tau+1})}\right)\norm{\sum^{t-1}_{\tau = t-\delta_t}\eta_{\tau}\textrm{lmo}(m^{\tau+1})}.
\end{eqnarray*}
%  
Assuming that stepsize $\eta_k$ is non-increasing, we have 
\begin{eqnarray*}
    \norm{\cT_2}_* &\le& \sum_{t=1}^k \prod_{j=t+1}^{k}(1-\alpha_j) \alpha_t \( L_0 + L_1 \norm{\nabla f(x^t)}_* \) \exp\left(L_1 R_t\eta_{t-\delta_t}\right)R_t\eta_{t-\delta_t} ~.
\end{eqnarray*}
% 
Bounding $\norm{\cT_3}_*$, we obtain 
\begin{eqnarray*}
    \norm{\cT_3}_* &\le& {\color{red} \rho} \norm{\sum_{t=1}^k \prod_{j=t+1}^{k}(1-\alpha_j) \alpha_t e_t}_2.
\end{eqnarray*}
% 
Taking expectation, we have 
\begin{eqnarray*}
    \Exp{\norm{\cT_3}_*} &\le& {\color{red} \rho} \Exp{\norm{\sum_{t=1}^k \prod_{j=t+1}^{k}(1-\alpha_j) \alpha_t e_t}_2}\le {\color{red} \rho} \sqrt{\Exp{\norm{\sum_{t=1}^k \prod_{j=t+1}^{k}(1-\alpha_j) \alpha_t e_t}_2^2}}\\
    &\le& {\color{red} \rho} \sqrt{\sum_{t=1}^k \prod_{j=t+1}^{k}(1-\alpha_j)^2 \alpha_t^2 \Exp{\norm{e_t}_2^2}}
        \le {\color{red} \rho} \sigma\sqrt{\sum_{t=1}^k \prod_{j=t+1}^{k}(1-\alpha_j)^2 \alpha_t^2} ~,
\end{eqnarray*}
where the second inequality uses Jensen's inequality and the third uses iid ness.
% 
Thus, we have 
\begin{eqnarray*}
    \Exp{\norm{\hat{e}_{k+1}}_*}
    &\le&  \sum_{t=1}^k \prod_{j=t+1}^{k}(1-\alpha_j) \left(L_0 + L_1\Exp{\norm{\nabla f(x^t)}_*}\right)\exp(L_1\eta_{t-1})\eta_{t-1}\\
    && + \sum_{t=1}^k \prod_{j=t+1}^{k}(1-\alpha_j) \alpha_t (L_0 + L_1 \Exp{\norm{\nabla f(x^t)}_*})\exp\left(L_1 R_t\eta_{t-\delta_t}\right)R_t\eta_{t-\delta_t}\\
    && + {\color{red} \rho} \sigma\sqrt{\sum_{t=1}^k \prod_{j=t+1}^{k}(1-\alpha_j)^2 \alpha_t^2} ~.
\end{eqnarray*}
% 
% 
\end{proof}
% 
% 
\section{Proof of \Cref{thm:iteration_complexity}}
\label{proof:iteration_complexity}
Let us first restate the theorem before proving it.
\begin{restate-theorem}{\ref{thm:iteration_complexity}}
    The iteration complexity is
    \begin{equation*}
        K = \cO\left(\frac{\left(\nicefrac{\Delta_0 }{\eta} + {\color{red} \rho} \sigma +  L_0\eta\right)^4 }{\varepsilon^4} \( \log \frac{1}{\varepsilon} \)^4\right).
     \end{equation*}
\end{restate-theorem}
% 
\subsection{Standard Smoothness Case}
% 
\begin{proof}
We set $L_1 = 0$, which implies standard smoothness of $f$.
\Cref{lem:descent_lemma} implies 
\begin{eqnarray*}
    \sum_{k=0}^{K-1} \eta_k \Exp{\norm{\nabla f(x^k)}_*} &\le& \Delta_0 + 2\sum^{K-1}_{k=0}\eta_k\Exp{\norm{\hat{e}_{k+1}}_*} +\frac{L_0}{2}\sum^{K-1}_{k=0}\eta_k^2
\end{eqnarray*}
% 
Then, by \Cref{lem:error_bound}, 
\begin{eqnarray*}
    \sum_{k=0}^{K-1} \eta_k \Exp{\norm{\nabla f(x^k)}_*} &\le& \Delta_0  +\frac{L_0}{2}\sum^{K-1}_{k=0}\eta_k^2 + 2\sum^{K-1}_{k=0}\eta_k\sum_{t=1}^k \prod_{j=t+1}^{k}(1-\alpha_j) L_0\eta_{t-1}\\
    && + 2\sum^{K-1}_{k=0}\eta_k\sum_{t=1}^k \prod_{j=t+1}^{k}(1-\alpha_j) \alpha_t L_0 R_t\eta_{t-\delta_t}\\
    && + 2{\color{red} \rho} \sigma\sum^{K-1}_{k=0}\eta_k\sqrt{\sum_{t=1}^k \prod_{j=t+1}^{k}(1-\alpha_j)^2 \alpha_t^2} ~.
\end{eqnarray*}
% 
Taking $R_t = \left\lfloor\frac{1}{\alpha_t}\right\rfloor$ and $\eta_k = \frac{\eta}{(k+1)^{\nicefrac{3}{4}}}$, we have 
\arto{Lemma 10 from (\href{https://arxiv.org/pdf/2311.03252}{Florian})}
\begin{eqnarray*}
    \sum_{t=1}^k \prod_{j=t+1}^{k}(1-\alpha_j)^2 \alpha_t^2 &\le& \sum_{t=1}^k \prod_{j=t+1}^{k}(1-\alpha_j) \alpha_t^2 =  \sum_{t=1}^k \prod_{j=t+1}^{k}(1-j^{-\nicefrac{1}{2}}) t^{-1}\\
        &\le& 2\exp\left(\frac{1}{1-\nicefrac{1}{2}}\right) \frac{1}{\sqrt{k+1}} \le \frac{18}{(k+1)^{\nicefrac{1}{2}}};\\
    \sum_{t=1}^k \prod_{j=t+1}^{k}(1-\alpha_j) \alpha_t R_t\eta_{t-\delta_t}
        &=& \sum_{t=1}^k \prod_{j=t+1}^{k}(1-\alpha_j) \alpha_t \left\lfloor\frac{1}{\alpha_t}\right\rfloor\eta_{t-\delta_t}\\
        &=& \sum_{t=1}^k \prod_{j=t+1}^{k}(1-j^{-\nicefrac{1}{2}}) \frac{\left\lfloor\sqrt{t}\right\rfloor}{\sqrt{t}} \eta (t-\delta_t+1)^{-\nicefrac{3}{4}}\\
        &\le& \eta \sum_{t=1}^k \prod_{j=t+1}^{k}(1-j^{-\nicefrac{1}{2}}) (t-\delta_t+1)^{-\nicefrac{3}{4}}\\
        &\le& \eta \sum_{t=1}^k \prod_{j=t+1}^{k}(1-j^{-\nicefrac{1}{2}}) (t-R_t+1)^{-\nicefrac{3}{4}}\\
        &\le& \eta \sum_{t=1}^k \prod_{j=t+1}^{k}(1-j^{-\nicefrac{1}{2}}) (t-\sqrt{t}+1)^{-\nicefrac{3}{4}} \\
        &\le& 2\eta \sum_{t=1}^k \prod_{j=t+1}^{k}(1-j^{-\nicefrac{1}{2}}) t^{-\nicefrac{3}{4}} \\
        &\le& 4\eta\exp\left(\frac{1}{1-\nicefrac{1}{2}}\right) (k+1)^{\nicefrac{1}{2} - \nicefrac{3}{4}}\\
        &\le& \frac{36\eta}{(k+1)^{\nicefrac{1}{4}}};\\
    \sum_{t=1}^k \prod_{j=t}^{k}(1-\alpha_j)\eta_{t-1} &=& \eta \sum_{t=1}^k \prod_{j=t+1}^{k}(1-j^{-\nicefrac{1}{2}})  t^{-\nicefrac{3}{4}}
        \le \frac{18\eta}{(k+1)^{\nicefrac{1}{4}}} ~.
\end{eqnarray*}
% 
Thus, we obtain
\begin{eqnarray*}
    \sum_{k=0}^{K-1} \eta_k \Exp{\norm{\nabla f(x^k)}_*} &\le& \Delta_0  +\frac{L_0}{2}\sum^{K-1}_{k=0}\eta_k^2 + 72L_0\eta\sum^{K-1}_{k=0}\frac{\eta_k}{(k+1)^{\nicefrac{1}{4}}}\\
    && + 36L_0\eta\sum^{K-1}_{k=0}\frac{\eta_k}{(k+1)^{\nicefrac{1}{4}}}\\
    && + 2{\color{red} \rho} \sigma\sum^{K-1}_{k=0}\eta_k\sqrt{\frac{18}{(k+1)^{\nicefrac{1}{2}}}}\\
    &\le& \Delta_0  +\frac{L_0\eta^2}{2}\sum^{K-1}_{k=0}\frac{1}{(k+1)^{\nicefrac{3}{2}}} + 98 L_0\eta^2 \sum^{K-1}_{k=0}\frac{1}{k+1} + 9 {\color{red} \rho} \sigma \eta \sum^{K-1}_{k=0}\frac{1}{k+1} ~.
\end{eqnarray*}
% 
Bounding sums as follows 
\begin{eqnarray*}
    \sum^{K-1}_{k = 0} \frac{1}{(k+1)^q} &=& \sum^{K}_{k = 1}  \frac{1}{k^q} \le \begin{cases}
        1 + \int\limits^{K}_1 \frac{dt}{t} = 1+  \log K, & \text{if } q =1 \ , \\ 
        1 + \int\limits^{K}_1 \frac{dt}{t^q} = 1+  \frac{1}{q-1} - \frac{1}{K^{q-1}} \leq \frac{q}{q-1}, & \text{if } q \geq 1 \ ,
    \end{cases}
\end{eqnarray*}
we have 
\begin{eqnarray*}
    \eta K^{\nicefrac{1}{4}}\min_{k \in \{0, \dots, K\}}\Exp{\norm{\nabla f(x^k)}_*} &\le& K\eta_K \min_{k \in \{0, \dots, K\}}\Exp{\norm{\nabla f(x^k)}_*} \le \sum_{k=0}^{K-1} \eta_k \Exp{\norm{\nabla f(x^k)}_*} \\
    &\le& \Delta_0   + 101 L_0\eta^2 (1+\log K) + 9 {\color{red} \rho} \sigma \eta (1+\log K) ,
\end{eqnarray*}
where $K\ge 2$.
% 
Thus, we have 
\begin{eqnarray*}
    \min_{k \in \{0, \dots, K\}}\Exp{\norm{\nabla f(x^k)}_*} 
    &\le& \frac{\nicefrac{\Delta_0 }{\eta}  + 101 L_0\eta (1+\log K) + 9 {\color{red} \rho} \sigma  (1+\log K)}{ K^{\nicefrac{1}{4}}},
\end{eqnarray*}
 and iteration complexity is defined as 
 \begin{equation*}
    K = \cO\left(\frac{\left(\nicefrac{\Delta_0 }{\eta} + {\color{red} \rho} \sigma +  L_0\eta\right)^4 }{\varepsilon^4} \( \log \frac{1}{\varepsilon} \)^4\right).
 \end{equation*}
% 
\end{proof}
% 
% 
\section{Time Complexity Analysis for the Fixed-Time Case}
\arto{name?}
% 
We now convert iteration complexity into wall-clock complexity under the fixed computation model from \eqref{eq:fixed_time}.
We begin with a basic bound for a fixed delay threshold.
% 
\begin{lemma}[Duration of $R$ updates {\citep[Lemma 4.1]{maranjyan2025ringmaster}}]
    \label{lem:time_R}
    Let $R$ be a fixed delay threshold.
    Under the fixed computation model \eqref{eq:fixed_time}, the time required to complete any $R$ consecutive updates is at most
    \begin{equation}\label{eq:time_R}
       t(R) \coloneqq 2 \min_{m \in [n]} \left(\frac{1}{m}\sum_{i=1}^{m}\frac{1}{\tau_i}\right)^{-1} \left( 1 + \frac{R}{m} \right).
    \end{equation}
\end{lemma}
% 
We next extend this bound to adaptive, nondecreasing thresholds.
% 
\begin{lemma}[Time for General Adaptive Threshold]\label{lem:time_RK}
    Let $R_k$ be a non-decreasing delay threshold.
    For each integer $r \in \{\lfloor R_0 \rfloor, \dots, \lfloor R_{K-1} \rfloor\}$, let $\mathcal{K}_r = \{k \mid \lfloor R_k \rfloor = r\}$ denote the block of iterations where the threshold is constant.
    Then the time required by \Cref{alg:ringmaster_lmo} to complete its first $K$ updates satisfies
    \begin{equation*}
        T(K) \le \sum_{r=\lfloor R_0 \rfloor}^{\lfloor R_{K-1} \rfloor} \left\lceil \frac{|\mathcal{K}_r|}{r} \right\rceil t(r) ~.
    \end{equation*}
\end{lemma}
% 
\begin{proof}
    Partition the updates $\{0,1,\dots,K-1\}$ into disjoint blocks $\mathcal{K}_r$ according to the integer value of $\lfloor R_k \rfloor$.
    Since $R_k$ is non-decreasing, each $\mathcal{K}_r$ is a contiguous block of iterations, and for every $k \in \mathcal{K}_r$ we have $r \le R_k < r+1$.

    By \Cref{lem:time_R}, the time needed to complete any $r$ consecutive updates with threshold $r$ is at most $t(r)$.
    The block $\mathcal{K}_r$ can be covered by
    \begin{equation*}
        N_b = \left\lceil \frac{|\mathcal{K}_r|}{r} \right\rceil
    \end{equation*}
    sub-blocks of size at most $r$.
    The last sub-block may contain fewer than $r$ updates, but its duration is still bounded by $t(r)$.
    Therefore, the total time spent on block $\mathcal{K}_r$ is at most $N_b t(r)$.
    Summing over all threshold values attained during the first $K$ updates yields
    \begin{equation*}
        T(K) = \sum_{r=\lfloor R_0 \rfloor}^{\lfloor R_{K-1} \rfloor} T(\mathcal{K}_r) \le \sum_{r=\lfloor R_0 \rfloor}^{\lfloor R_{K-1} \rfloor} \left\lceil \frac{|\mathcal{K}_r|}{r} \right\rceil t(r) ~.
    \end{equation*}
\end{proof}

\subsection{Proof of \Cref{lemma:time_sqrt}}\label{proof:sqrt_threshold}
% 
We now specialize \Cref{lem:time_RK} to the square-root threshold schedule.
\begin{restate-lemma}{\ref{lemma:time_sqrt}}[Time Complexity for Square-Root Threshold]
    Let the delay threshold be $R_k = \lfloor \sqrt{k+1} \rfloor$ for $k \ge 0$.
    Under the fixed computation model \eqref{eq:fixed_time}, the wall-clock time needed to complete $K$ iterations satisfies
    \begin{equation*}
        T(K) = \mathcal{O}\left( \min_{m \in [n]} \left(\frac{1}{m}\sum_{i=1}^{m}\frac{1}{\tau_i}\right)^{-1} \left( \sqrt{K} + \frac{K}{m} \right) \right).
    \end{equation*}
\end{restate-lemma}
% 
\begin{proof}
    Apply \Cref{lem:time_RK} with $R_k = \lfloor \sqrt{k+1} \rfloor$.
    For a fixed integer $r \ge 1$, the condition $\lfloor \sqrt{k+1} \rfloor = r$ is equivalent to
    $r \le \sqrt{k+1} < r+1$, or equivalently $r^2 \le k+1 < (r+1)^2$.
    Hence,
    \begin{equation*}
        k \in \{r^2-1, r^2, \dots, (r+1)^2-2\}.
    \end{equation*}
    Therefore, the block size is
    \begin{equation*}
        |\mathcal{K}_r| = \left( (r+1)^2 - 2 \right) - (r^2 - 1) + 1 = (r+1)^2 - r^2 = 2r + 1.
    \end{equation*}
    For $r \ge 1$, the ceiling factor in \Cref{lem:time_RK} satisfies
    \begin{equation*}
        \left\lceil \frac{2r+1}{r} \right\rceil = \left\lceil 2 + \frac{1}{r} \right\rceil = 3.
    \end{equation*}
    Let $r_{\max} = \lfloor \sqrt{K} \rfloor$.
    Using \Cref{lem:time_RK} and \Cref{eq:time_R}, we obtain
    \begin{equation*}
        T(K) \le \sum_{r=1}^{r_{\max}} 3\, t(r)
        = 6 \sum_{r=1}^{r_{\max}} \min_{m \in [n]} \left[ \left(\frac{1}{m}\sum_{i=1}^{m}\frac{1}{\tau_i}\right)^{-1} \left( 1 + \frac{r}{m} \right) \right].
    \end{equation*}
    Since $\sum_r \min_m a_{r,m} \le \min_m \sum_r a_{r,m}$, it follows that
    \begin{align*}
        T(K)
        &\le 6 \min_{m \in [n]} \left[
            \sum_{r=1}^{r_{\max}} \left(\frac{1}{m}\sum_{i=1}^{m}\frac{1}{\tau_i}\right)^{-1}
            + \sum_{r=1}^{r_{\max}} \left(\frac{1}{m}\sum_{i=1}^{m}\frac{1}{\tau_i}\right)^{-1} \frac{r}{m}
        \right] \\
        &= 6 \min_{m \in [n]} \left[
            \left(\frac{1}{m}\sum_{i=1}^{m}\frac{1}{\tau_i}\right)^{-1} r_{\max}
            + \left(\frac{1}{m}\sum_{i=1}^{m}\frac{1}{\tau_i}\right)^{-1} \frac{r_{\max}(r_{\max}+1)}{2m}
        \right] \\
        &\le 6 \min_{m \in [n]} \left[
            \left(\frac{1}{m}\sum_{i=1}^{m}\frac{1}{\tau_i}\right)^{-1} \sqrt{K}
            + \left(\frac{1}{m}\sum_{i=1}^{m}\frac{1}{\tau_i}\right)^{-1} \frac{K}{m}
        \right].
    \end{align*}
    This proves that
    \begin{equation*}
        T(K) = \mathcal{O}\left( \min_{m \in [n]} \left(\frac{1}{m}\sum_{i=1}^{m}\frac{1}{\tau_i}\right)^{-1} \left( \sqrt{K} + \frac{K}{m} \right) \right).
    \end{equation*}
\end{proof}
% 
% 
% 
\section{Time Complexity Analysis for Arbitrarily Changing Computation Times}\label{sec:arbitrary_time}
% 
In practice, the \emph{fixed computation model} \eqref{eq:fixed_time} is often not satisfied.
The compute power of devices can vary over time due to temporary disconnections, hardware or network delays, fluctuations in processing capacity, or other transient effects \citep{maranjyan2025mindflayer}.

In this section we extend our theory to the more general setting of arbitrarily varying computation times.

\subsection{Universal Computation Model}
\label{sec:universal_computation_model}
To formalize this setting, we adopt the \emph{universal computation model} introduced by \citet{tyurin2024tighttimecomplexitiesparallel}.

For each worker $i \in [n]$, we define a \emph{compute power} function
$$
    p_i : \R_{+} \to \R_{+}~,
$$
assumed nonnegative and continuous almost everywhere (countably many jumps allowed).
For any $T_2 \ge T_1 \ge 0$, the number of stochastic gradients \emph{completed} by worker $i$ on $[T_1, T_2]$ is
$$
    \#\text{gradients in }[T_1, T_2] \;=\; \left\lfloor \int_{T_1}^{T_2} p_i(t)\,dt \right\rfloor.
$$
Here, $p_i(t)$ models the worker's time-varying computational ability: smaller values over an interval yield fewer completed gradients, and larger values yield more.

For instance, if worker $i$ remains idle for the first $T$ seconds and then becomes active, this corresponds to $p_i(t) = 0$ for $t \leq T$ and $p_i(t) > 0$ for $t > T$.
More generally, $p_i(t)$ may follow periodic or irregular patterns, leading to bursts of activity, pauses, or chaotic changes in compute power.
The process $p_i(t)$ may even be random, and all results hold conditional on the realized sample paths of $\{p_i\}$.

The \emph{universal computation model} reduces to the \emph{fixed computation model} \eqref{eq:fixed_time} when $p_i(t) = \nicefrac{1}{\tau_i}$ for all $t \geq 0$ and $i \in [n]$.
In this case,
$$
    \#\text{gradients in }[T_1, T_2] = \left\lfloor \frac{T_2 - T_1}{\tau_i} \right\rfloor,
$$  
meaning that worker $i$ computes one stochastic gradient after $T_1 + \tau_i$ seconds, two gradients after $T_1 + 2\tau_i$ seconds, and so on.
\arto{rewrite all above}
% 
\subsection{Analysis}
% 
\begin{lemma}[Duration of $R$ updates {\citep[Lemma 5.1]{maranjyan2025ringmaster}}]
    \label{lem:time_R_arbitrary}
    Let the workers' computation times satisfy the \emph{universal computation model}.
    Let $R$ be the delay threshold of \Cref{alg:ringmaster_lmo}.
    Assume that some iteration starts at time $T_0$.
    Starting from this iteration, the $R$ consecutive iterate updates of \Cref{alg:ringmaster_lmo} will be performed within the elapsed time
    \begin{equation*} 
           t(R;T_0) \eqdef \min \left\{t \geq 0 : \sum \limits_{i=1}^{n} \flr{\frac{1}{4} \int_{T_0}^{T_0 + t} p_i(\tau)\, d \tau} \geq R\right\}.
    \end{equation*}
\end{lemma}

Then, under the \emph{universal computation model}, \ringmaster takes $T_{K}$ seconds for $K$ iterations, where $T_{K}$ is the $K$-th element of the following recursively defined sequence:
\begin{align*}
    T_{k} \eqdef \min \left\{T \geq 0 : \sum \limits_{i=1}^{n} \flr{\frac{1}{4} \int_{T_{k - 1}}^{T} p_i(\tau) d \tau} \geq R\right\} .
\end{align*}
for all $k \geq 1$ and $T_{0} = 0$.
This result is from \citep{maranjyan2025ringmaster}.


Now we want to find the time for changig threshold $R_k$ case.

We next extend \Cref{lem:time_R_arbitrary} to non-decreasing adaptive thresholds.

\begin{lemma}[Time for General Adaptive Threshold under the Universal Computation Model]
    \label{lem:time_RK_arbitrary}
    Let $R_k$ be a non-decreasing delay threshold.
    For each integer $r \in \{\lfloor R_0 \rfloor, \dots, \lfloor R_{K-1} \rfloor\}$, let
    $$
        \mathcal{K}_r \coloneqq \{k \mid \lfloor R_k \rfloor = r\}
    $$
    denote the block of iterations where the threshold is constant.
    For every such $r$, partition $\mathcal{K}_r$ into
    $$
        N_r \coloneqq \left\lceil \frac{|\mathcal{K}_r|}{r} \right\rceil
    $$
    consecutive sub-blocks, each containing at most $r$ iterations.

    Let $B \coloneqq \sum_{r=\lfloor R_0 \rfloor}^{\lfloor R_{K-1} \rfloor} N_r$, and enumerate all these sub-blocks in chronological order.
    Denote by $r_j$ the threshold value attached to the $j$-th sub-block, where $j = 1, \dots, B$.
    Define the sequence $\{T_j\}_{j=0}^{B}$ recursively by
    \begin{align*}
        T_0 &\coloneqq 0, \\
        T_j &\coloneqq T_{j-1} + t(r_j; T_{j-1}), \qquad j = 1, \dots, B.
    \end{align*}
    Then the wall-clock time needed to complete the first $K$ iterations of \Cref{alg:ringmaster_lmo} satisfies
    \begin{equation*}
        T(K) \le T_B.
    \end{equation*}
\end{lemma}

\begin{proof}
    Partition the iterations $\{0,1,\dots,K-1\}$ into disjoint blocks $\mathcal{K}_r$ according to the integer value of $\lfloor R_k \rfloor$.
    Since $R_k$ is non-decreasing, each $\mathcal{K}_r$ is a contiguous block of iterations.

    Fix some $r$ and consider one of its sub-blocks.
    By construction, this sub-block contains at most $r$ consecutive iterations, and throughout this sub-block the threshold satisfies
    $$
        r \le R_k < r+1.
    $$
    Therefore, using threshold $r$ can only make the algorithm more conservative than the actual threshold schedule on this sub-block.
    Hence \Cref{lem:time_R_arbitrary} implies that, if this sub-block starts at time $T$, then all its iterations are completed within time at most $t(r;T)$.

    Now enumerate all sub-blocks in chronological order and let $r_j$ be the threshold associated with the $j$-th one.
    Starting from time $T_0 = 0$, the first sub-block is completed by time at most
    $$
        T_1 = T_0 + t(r_1;T_0).
    $$
    Repeating the same argument inductively, if the first $j-1$ sub-blocks are completed by time at most $T_{j-1}$, then the $j$-th sub-block is completed by time at most
    $$
        T_j = T_{j-1} + t(r_j;T_{j-1}).
    $$
    After all
    $$
        B = \sum_{r=\lfloor R_0 \rfloor}^{\lfloor R_{K-1} \rfloor} \left\lceil \frac{|\mathcal{K}_r|}{r} \right\rceil
    $$
    sub-blocks are processed, all first $K$ iterations have been completed. Therefore,
    \begin{equation*}
        T(K) \le T_B.
    \end{equation*}
\end{proof}

This is the arbitrary-time analogue of \Cref{lem:time_RK} from the fixed computation model.
It reduces the analysis of a changing threshold schedule to a recursion over threshold-constant sub-blocks, where each sub-block is controlled by the fixed-threshold bound \Cref{lem:time_R_arbitrary}.

\begin{lemma}[Square-Root Threshold under the Universal Computation Model]
    \label{lem:time_sqrt_arbitrary}
    Let $R_k = \lfloor \sqrt{k+1} \rfloor$ for $k \ge 0$.
    Under the universal computation model, the wall-clock time needed to complete the first $K$ iterations of \Cref{alg:ringmaster_lmo} is bounded by
    \begin{equation*}
        T(K) \le S_{3\lfloor \sqrt{K} \rfloor},
    \end{equation*}
    where the sequence $\{S_j\}_{j=0}^{3\lfloor \sqrt{K} \rfloor}$ is defined recursively as
    \begin{align*}
        S_0 &\coloneqq 0, \\
        S_j &\coloneqq S_{j-1} + t\!\left(\left\lceil \frac{j}{3} \right\rceil; S_{j-1}\right), \qquad j = 1, \dots, 3\lfloor \sqrt{K} \rfloor .
    \end{align*}
\end{lemma}

\begin{proof}
    We apply \Cref{lem:time_RK_arbitrary} with $R_k = \lfloor \sqrt{k+1} \rfloor$.
    For every integer $r \ge 1$, the condition $\lfloor \sqrt{k+1} \rfloor = r$ is equivalent to
    $$
        r^2 \le k+1 < (r+1)^2,
    $$
    hence
    $$
        \mathcal{K}_r = \{r^2-1, r^2, \dots, (r+1)^2-2\},
    $$
    and therefore
    \begin{equation*}
        |\mathcal{K}_r| = (r+1)^2-r^2 = 2r+1.
    \end{equation*}
    Consequently,
    \begin{equation*}
        \left\lceil \frac{|\mathcal{K}_r|}{r} \right\rceil
        =
        \left\lceil \frac{2r+1}{r} \right\rceil
        =
        3.
    \end{equation*}
    Thus each threshold level $r$ contributes exactly three sub-blocks of size at most $r$.

    Let $r_{\max} \coloneqq \lfloor \sqrt{K} \rfloor$.
    Enumerating all sub-blocks in chronological order gives exactly $3r_{\max}$ sub-blocks, with threshold values
    $$
        1,1,1,2,2,2,\dots,r_{\max},r_{\max},r_{\max}.
    $$
    Therefore, the recursion from \Cref{lem:time_RK_arbitrary} becomes
    \begin{align*}
        S_0 &\coloneqq 0, \\
        S_j &\coloneqq S_{j-1} + t\!\left(\left\lceil \frac{j}{3} \right\rceil; S_{j-1}\right), \qquad j = 1, \dots, 3r_{\max},
    \end{align*}
    and the resulting completion time satisfies
    \begin{equation*}
        T(K) \le S_{3r_{\max}} = S_{3\lfloor \sqrt{K} \rfloor}.
    \end{equation*}
\end{proof}

\begin{theorem}[Final Arbitrary-Time Wall-Clock Complexity]
    \label{thm:time_arbitrary}
    Consider \Cref{alg:ringmaster_lmo} with
    $$
        \alpha_0 = 1,\qquad
        \alpha_k = \frac{1}{\sqrt{k}},\qquad
        R_k = \lfloor \sqrt{k+1} \rfloor,
        \qquad
        \eta_k = \frac{\eta}{(k+1)^{3/4}}
    $$
    for $k \ge 1$.
    Let
    \begin{equation*}
        K_{\varepsilon}
        =
        \widetilde{\mathcal{O}}\left(\frac{\Phi^4}{\varepsilon^4}\right)
    \end{equation*}
    be the iteration complexity from \Cref{thm:iteration_complexity}.
    Then, under the universal computation model, the wall-clock time needed to reach an $\varepsilon$-stationary point satisfies
    \begin{equation*}
        T_{\varepsilon}
        \le
        S_{3\lfloor \sqrt{K_{\varepsilon}} \rfloor}
        =
        S_{\widetilde{\mathcal{O}}(\Phi^2/\varepsilon^2)},
    \end{equation*}
    where the sequence $\{S_j\}$ is defined recursively by
    \begin{align*}
        S_0 &\coloneqq 0, \\
        S_j &\coloneqq S_{j-1} + t\!\left(\left\lceil \frac{j}{3} \right\rceil; S_{j-1}\right).
    \end{align*}
\end{theorem}

\begin{proof}
    By \Cref{thm:iteration_complexity}, the method reaches an $\varepsilon$-stationary point within
    \begin{equation*}
        K_{\varepsilon}
        =
        \widetilde{\mathcal{O}}\left(\frac{\Phi^4}{\varepsilon^4}\right)
    \end{equation*}
    iterations.
    Applying \Cref{lem:time_sqrt_arbitrary} with $K = K_{\varepsilon}$ yields
    \begin{equation*}
        T_{\varepsilon}
        \le
        S_{3\lfloor \sqrt{K_{\varepsilon}} \rfloor}.
    \end{equation*}
    Since $\sqrt{K_{\varepsilon}} = \widetilde{\mathcal{O}}(\Phi^2/\varepsilon^2)$, the claim follows.
\end{proof}

\paragraph{Comparison with \ringmaster.}
The fixed-threshold arbitrary-time result from \citet{maranjyan2025ringmaster} controls progress through blocks of a single threshold $R$, leading to a recursion over roughly $K/R$ blocks of size $R$.
In contrast, with the square-root schedule $R_k = \lfloor \sqrt{k+1} \rfloor$, \Cref{lem:time_sqrt_arbitrary} shows that the first $K$ iterations are covered by only
\begin{equation*}
    3\lfloor \sqrt{K} \rfloor
\end{equation*}
adaptive sub-blocks.
Thus our schedule replaces the fixed-threshold decomposition by an adaptive one whose depth grows only as $\mathcal{O}(\sqrt{K})$.
After substituting the iteration complexity bound, this yields an arbitrary-time wall-clock recursion of length
\begin{equation*}
    \widetilde{\mathcal{O}}\left(\frac{\Phi^2}{\varepsilon^2}\right),
\end{equation*}
whereas the fixed-threshold \ringmaster analysis scales with the number of fixed-$R$ blocks.
This captures the advantage of allowing the delay threshold to increase with $k$: the method adapts to larger delays over time without requiring a single globally tuned threshold.


\end{document}
